By Lemma~\texttt{lem:beta-class-one-sided-monotonicity}, the order-\(\beta^{\prime}\) class is contained in the order-\(\beta\) class. Therefore, for every \(n\),
\[
\mathcal M_n(\alpha,\beta^{\prime},s,d,M_Y,L,c,t_0)
\le
\mathcal M_n(\alpha,\beta,s,d,M_Y,L,c,t_0).
\]
The hypothesis \((\alpha,\beta,s,d)\in\Omega\) permits application of Theorem~\texttt{thm:oracle-high-s} to the rougher order-\(\beta\) class. Hence there are an \(n\)-independent constant \(C_{\mathrm{up}}^0<\infty\) and \(n_0\in\mathbb N\) such that, for every \(n\ge n_0\),
\[
\mathcal M_n(\alpha,\beta,s,d,M_Y,L,c,t_0)
\le C_{\mathrm{up}}^0r_n(\alpha).
\]
Independently, Lemma~\texttt{lem:all-beta-oracle-lower-floor}, applied at the smoother order \(\beta^{\prime}\), supplies an \(n\)-independent constant \(C_{\mathrm{low}}>0\) such that, for every \(n\ge1\),
\[
\mathcal M_n(\alpha,\beta^{\prime},s,d,M_Y,L,c,t_0)
\ge C_{\mathrm{low}}r_n(\alpha).
\]
Set \(C_{\mathrm{up}}=\max\{C_{\mathrm{up}}^0,C_{\mathrm{low}}\}\). Substitution of \(r_n(\alpha)=n^{-2\alpha/(2\alpha+1)}\) from Definition~\texttt{def:oracle-rate} yields, for every \(n\ge n_0\),
\[
C_{\mathrm{low}}n^{-2\alpha/(2\alpha+1)}
\le \mathcal M_n(\alpha,\beta^{\prime},s,d,M_Y,L,c,t_0)
\le \mathcal M_n(\alpha,\beta,s,d,M_Y,L,c,t_0)
\le C_{\mathrm{up}}n^{-2\alpha/(2\alpha+1)}.
\]
The smoother-class risk is positive by the first inequality. Dividing through by it gives
\[
1\le
\frac{\mathcal M_n(\alpha,\beta,s,d,M_Y,L,c,t_0)}
{\mathcal M_n(\alpha,\beta^{\prime},s,d,M_Y,L,c,t_0)}
\le\frac{C_{\mathrm{up}}}{C_{\mathrm{low}}}.
\]
The oracle upper theorem was used only at the rough order \(\beta\); the smoother order \(\beta^{\prime}\) entered only through the all-order lower floor. Thus no oracle-region condition at \(\beta^{\prime}\) is needed.